% =====================================================================
% PI-NCA: which structural prior matches which PDE regime?
%
% Every number in this file comes from a macro or a table generated by
%   python -m pinca_jax.paper
% from results/*.json. Nothing quantitative is typed by hand. If a number looks
% stale, regenerate; if a macro is undefined, the experiment has not been run.
%
% Build:
%   python -m pinca_jax.paper && pdflatex main && bibtex main && pdflatex main
% Drop iclr2026_conference.sty next to this file for the conference format; the
% document falls back to article class if it is absent.
% =====================================================================
\documentclass{article}

\IfFileExists{iclr2026_conference.sty}{%
  \usepackage{iclr2026_conference,times}%
}{%
  \usepackage[margin=1in]{geometry}%
  \typeout{NOTE: iclr2026_conference.sty not found; building in article class.}%
}

\usepackage{amsmath,amssymb}
\usepackage{booktabs}
\usepackage{graphicx}
\usepackage{xcolor}
\usepackage{url}
\usepackage[hidelinks]{hyperref}
\usepackage{caption}

% Generated by `python -m pinca_jax.paper`. Do not edit.
% Every number quoted in the paper prose is defined here, from results/.
\newcommand{\Backends}{cpu}
\newcommand{\Grids}{8, 12, 16, 24}
\newcommand{\MaxSeeds}{3}
\newcommand{\MinSeeds}{1}
\newcommand{\Nablations}{5}
\newcommand{\Narchs}{14}
\newcommand{\Ncitations}{52}
\newcommand{\Ncitationsresolved}{52}
\newcommand{\Nfailedcells}{0}
\newcommand{\NpdesThreeD}{6}
\newcommand{\NpdesTwoD}{10}
\newcommand{\Nresultfiles}{52}
\newcommand{\Ntests}{81}
\newcommand{\beatsFloorAdvDiff}{yes}
\newcommand{\beatsFloorAllenCahn}{yes}
\newcommand{\beatsFloorCahnHilliard}{yes}
\newcommand{\beatsFloorFitzhughNagumo}{yes}
\newcommand{\beatsFloorGrayScott}{yes}
\newcommand{\beatsFloorHeat}{yes}
\newcommand{\beatsFloorNagumo}{yes}
\newcommand{\beatsFloorNavierStokes}{yes}
\newcommand{\beatsFloorShallowWater}{yes}
\newcommand{\beatsFloorWave}{yes}
\newcommand{\bestAdvDiff}{FNO}
\newcommand{\bestAdvDiffCI}{--}
\newcommand{\bestAdvDiffErr}{0.00657}
\newcommand{\bestAllenCahn}{FNO}
\newcommand{\bestAllenCahnCI}{--}
\newcommand{\bestAllenCahnErr}{0.00685}
\newcommand{\bestCahnHilliard}{ResNet}
\newcommand{\bestCahnHilliardCI}{[0.041, 0.046]}
\newcommand{\bestCahnHilliardErr}{0.0435}
\newcommand{\bestFitzhughNagumo}{NCA}
\newcommand{\bestFitzhughNagumoCI}{--}
\newcommand{\bestFitzhughNagumoErr}{0.125}
\newcommand{\bestGrayScott}{ResNet (iso-param)}
\newcommand{\bestGrayScottCI}{[0.2, 0.22]}
\newcommand{\bestGrayScottErr}{0.214}
\newcommand{\bestHeat}{Spectral PI-NCA}
\newcommand{\bestHeatCI}{[0.039, 0.044]}
\newcommand{\bestHeatErr}{0.0417}
\newcommand{\bestNagumo}{ResNet (iso-param)}
\newcommand{\bestNagumoCI}{[0.02, 0.024]}
\newcommand{\bestNagumoErr}{0.0221}
\newcommand{\bestNavierStokes}{U-Net}
\newcommand{\bestNavierStokesCI}{[0.094, 0.14]}
\newcommand{\bestNavierStokesErr}{0.115}
\newcommand{\bestShallowWater}{PI-NCA (flux)}
\newcommand{\bestShallowWaterCI}{[0.01, 0.014]}
\newcommand{\bestShallowWaterErr}{0.0122}
\newcommand{\bestWave}{PI-NCA (flux)}
\newcommand{\bestWaveCI}{[0.032, 0.044]}
\newcommand{\bestWaveErr}{0.0371}
\newcommand{\floorAdvDiff}{0.137}
\newcommand{\floorAllenCahn}{0.0266}
\newcommand{\floorCahnHilliard}{0.412}
\newcommand{\floorFitzhughNagumo}{1.05}
\newcommand{\floorGrayScott}{0.391}
\newcommand{\floorHeat}{0.364}
\newcommand{\floorNagumo}{0.0797}
\newcommand{\floorNavierStokes}{0.156}
\newcommand{\floorShallowWater}{0.0283}
\newcommand{\floorWave}{0.0755}
\newcommand{\matchedCrossover}{1}
\newcommand{\matchedEmuErr}{0.161}
\newcommand{\matchedK}{4}
\newcommand{\matchedPinnErr}{0.318}
\newcommand{\matchedSolverSpeedup}{1541}
\newcommand{\scalingNunstable}{1}
\newcommand{\scalingUnstableAxes}{grid}
\newcommand{\stabBoundedFail}{0}
\newcommand{\stabHorizon}{96}
\newcommand{\stabNCAFail}{100}
\newcommand{\teacherHeatSpatial}{0.02}
\newcommand{\teacherHeatTemporal}{0.0072}
\newcommand{\teacherHeatTotal}{0.013}


\newcommand{\code}[1]{\texttt{\small #1}}
\newcommand{\limitation}[1]{\par\smallskip\noindent\textit{Limitation.}~#1\par\smallskip}

\title{Which structural prior matches which PDE regime?\\
       A controlled comparison of cellular-automaton, spectral and\\
       physics-free neural surrogates}

\author{%
  Tharun N.\ Nanju \\
  \texttt{tharunnnanju@gmail.com}
}

\date{}

\begin{document}
\maketitle

% =====================================================================
\begin{abstract}
Neural cellular automata (NCAs), neural operators and ordinary convolutional networks
are all used as autoregressive surrogates for partial differential equations, and each
carries a different structural prior: strict locality with a single shared update rule,
global spectral mixing, or none at all. We ask which prior pays off in which regime,
under one harness, with the controls that make a negative answer reportable. We compare
\Narchs{} architectures on \NpdesTwoD{} two-dimensional and \NpdesThreeD{}
three-dimensional PDE families, distilled from differentiable reference solvers, and
report every comparison as a paired test on shared initial conditions with bootstrap
confidence intervals and Holm--Bonferroni correction. Three controls are included by
construction: a do-nothing identity baseline, physics-free CNN and U-Net surrogates at
full and parameter-matched size, and the numerical solver itself with its own
discretisation error measured against closed-form solutions. We find no universal
winner, and we find that several apparent wins do not survive the controls. We also
report a defect in the standard recipe for combining boundedness with conservation---the
uniform mass re-projection applied after a clip pushes cells back outside the bound it
just enforced---and give a headroom-proportional projection that restores mass exactly
while remaining inside the bound. Long-horizon stress tests are run with the divergence
guard disabled, so failures are counted rather than clamped: on Cahn--Hilliard at
\stabHorizon{} steps the unconstrained NCA fails on \stabNCAFail\% of initial conditions
while the bounded conservative variant fails on \stabBoundedFail\%. All code, raw
results, and a machine-checked audit of every claim in this paper are released.
\end{abstract}

% =====================================================================
\section{Introduction}

A neural surrogate for a PDE is a bet on a structural prior. A neural cellular automaton
bets that a single homogeneous local update rule, applied identically at every cell, is
the right hypothesis class, because that is what a differential operator is. A Fourier
neural operator bets on global spectral mixing. A conservative flux-divergence update
bets that mass conservation should be exact by construction rather than learned. An
ordinary residual CNN bets on nothing in particular.

The literature evaluates these bets one at a time, usually in papers whose proposed
method is one of the options. That produces a collection of individually plausible wins
that do not compose into an answer to the question a practitioner actually has, which is
not ``is this architecture good'' but ``given this equation, which prior should I
choose''.

\paragraph{The falsifiable hypothesis.} We state it so it can fail:

\begin{quote}
\textit{The best-performing structural prior is determined by the equation's structure,
not by the architecture's capacity. Specifically: (i) where the dynamics are local and
the quantity conserved, a local conservative update attains lower error than a global
spectral operator at a fraction of the parameters; (ii) where the solution carries
fine-scale or globally coupled structure, global spectral mixing attains lower error;
(iii) where the dynamics are non-conservative, imposing conservation costs accuracy.}
\end{quote}

Each clause is falsified by a paired test on shared initial conditions showing the
opposite ordering, or showing no resolvable difference. We report both outcomes.

\paragraph{The decision rule the hypothesis implies.} If it holds, a practitioner can
choose without a sweep: check whether the quantity is conserved, whether the field is
bounded, how far information must propagate per step, and how much fine-scale content
the solution carries. Table~\ref{tab:regime} is that rule, populated by measurement.

\paragraph{Why the controls matter more than the architecture.} An architecture
comparison is only as good as what it is compared against, and three baselines are
missing from most such comparisons. We include all three, and each one changes a
conclusion.

\begin{itemize}
\item \textbf{A do-nothing identity floor}, $g(x)=x$. On phenomena whose solution changes
  slowly relative to the evaluation horizon, this is a strong baseline, and a model above
  it has learned something worse than nothing. Without the floor in the table, that
  outcome is invisible.
\item \textbf{Physics-free CNN and U-Net surrogates}, at full size and at parameter counts
  matched to the NCAs. Without them, ``the conservation prior helps'' is confounded with
  ``an ordinary CNN was never tried''.
\item \textbf{The numerical solver}, with its own error measured against closed-form
  solutions where they exist. A surrogate that is less accurate \emph{and} slower than the
  solver it imitates has no use case, and at the grid sizes used here that is frequently
  the situation (Section~\ref{sec:worth}).
\end{itemize}

\paragraph{Contributions.}
\begin{enumerate}
\item \textbf{A regime map} over \NpdesTwoD{} 2-D and \NpdesThreeD{} 3-D PDE families and
  \Narchs{} architectures under one distillation harness, with every comparison reported
  as a paired Wilcoxon test on shared initial conditions, bootstrap confidence intervals,
  and Holm--Bonferroni correction across the architectures compared per equation
  (Section~\ref{sec:regime}).
\item \textbf{A defect and a fix in bounded conservative updates.} Clipping a field to its
  physical range and then restoring mass with a uniform offset---the standard
  composition---re-violates the bound the clip just enforced. We give a
  headroom-proportional projection that is mass-exact, bound-preserving, differentiable,
  and a single vectorised pass, and we ablate it against the uniform choice
  (Section~\ref{sec:proj}).
\item \textbf{Failure counting instead of failure clamping.} Long-horizon evaluation is
  run with the stabilising output guard disabled, so divergence is reported as a failure
  rate curve rather than absorbed into a saturated field and a merely-poor score
  (Section~\ref{sec:stability}).
\item \textbf{An accounting of what the comparison is worth.} We measure the teacher's own
  discretisation error against exact solutions, separate its spatial and temporal parts,
  and report the amortisation crossover at which a trained emulator becomes cheaper than
  per-instance PINN solving---and note where neither beats simply running the solver
  (Sections~\ref{sec:teacher} and \ref{sec:worth}).
\item \textbf{A machine-checked release.} Every count and headline claim in this paper is
  regenerated from the raw results by \code{pinca\_jax.claims}; every table is generated by
  \code{pinca\_jax.paper}; every one of the \Ncitations{} cited arXiv identifiers is
  verified against the arXiv API by \code{pinca\_jax.bib} (\Ncitationsresolved{} resolve);
  and \Ntests{} correctness tests assert the ported solvers equal their references before
  any experiment runs.
\end{enumerate}

% =====================================================================
\section{Related work, and what is new here}

Detailed positioning, including a sixteen-row comparison table over locality,
conservation mechanism, supervision, suite and dimensionality, is in
Appendix~\ref{app:related} and \code{docs/related\_work.md}. The short version, stated
against the closest prior work rather than around it:

\citet{saha2026nca} propose an NCA surrogate for PDE time-stepping with the same
motivation used here---a small local homogeneous rule mirroring the locality of
differential operators---benchmarked against PDE-Net, a PINN and an FNO on five canonical
PDEs beyond the training horizon. \citet{richardson2023nca} train NCAs on PDE
trajectories, show generalisation beyond the training data, and constrain NCAs to respect
symmetries. So \emph{NCAs as viable PDE surrogates}, and \emph{NCAs versus PINNs and
FNOs}, are established, not contributions of this work.

Conservative structure in a learned PDE model is likewise established.
\citet{praditia2021finn} embed the finite-volume structure so each quantity follows its
own conservation law with explicit inter-cell fluxes; \citet{richterpowell2022ncl}
parameterise networks that satisfy the continuity equation exactly by construction. This
paper does not claim to introduce conservative neural PDE models.

What remains is narrower. First, the conservation prior \emph{inside} a single homogeneous
CA rule---a strictly stronger constraint than FINN's per-cell parameters---and a
controlled measurement of whether it still pays once that constraint is imposed (ablation
A4, Section~\ref{sec:ablate}). Second, the bounded-and-conserving tension and its
resolution (Section~\ref{sec:proj}), which we have not found addressed elsewhere. Third,
a regime characterisation with the identity floor, physics-free controls and paired
statistics that make a null result reportable, over a suite that includes multi-field
systems the five-PDE comparisons do not cover. This is a benchmark-and-analysis
contribution, and we frame it as one.

% =====================================================================
\section{Setup}
\label{sec:setup}

\paragraph{Task.} Every architecture is treated as a one-step autoregressive emulator
$g_\theta: u_t \mapsto u_{t+1}$ on a periodic grid, trained by distillation against a
differentiable reference solver through a rollout of $H$ steps with backpropagation
through the whole rollout, and evaluated by rolling out $K \gg H$ steps from held-out
initial conditions. PINNs do not fit this interface---they solve a single initial-value
problem in continuous $(x,t)$---and are therefore compared separately and explicitly,
under a matched task (Section~\ref{sec:worth}), rather than placed in the same accuracy
column.

\paragraph{Equations.} \NpdesTwoD{} 2-D families spanning the axes the hypothesis names:
conservative and non-conservative, bounded and unbounded, scalar and multi-field, local
and globally coupled, smooth and sharp-interfaced. \NpdesThreeD{} of them are repeated in
3-D. Full specifications, parameters and stability limits are in
Appendix~\ref{app:suite}.

\paragraph{Architectures.} \Narchs{} models in three families---unconstrained local
(NCA), conservative local (flux-divergence PI-NCA and its bounded, multi-scale and
spectral hybrids), and global (FNO)---plus the physics-free controls (ResNet, U-Net, each
also at a parameter count matched to the NCAs) and the identity floor. Update equations
are in Appendix~\ref{app:arch}. All architectures are channel-generic, so the comparison
matrix is rectangular: the same model list is measured on every phenomenon.

\paragraph{Statistics.} Each cell is trained over multiple seeds; each trained model is
evaluated on a fixed batch of held-out initial conditions shared by every architecture, in
the same order. Reported errors are per-initial-condition relative $L_2$, summarised with
a 10{,}000-sample percentile bootstrap. Architecture comparisons are paired Wilcoxon
signed-rank tests on those shared initial conditions, with Holm--Bonferroni correction
across the family of comparisons made per equation. A difference is called significant
only when the bootstrap interval on the paired differences excludes zero \emph{and}
Wilcoxon rejects; otherwise the entry reads \emph{tie}, which means unresolvable at this
sample size and not that the means are equal.

\paragraph{Compute and provenance.} Every results file records the JAX version, backend,
devices, peak memory, and the exact configuration. Current inventory: backend
\Backends{}, grids \Grids{}, seeds per table \MinSeeds--\MaxSeeds{},
\Nresultfiles{} result files, \Nfailedcells{} failed matrix cells.
\code{docs/claims\_audit.md} lists which pre-registered claims this inventory currently
supports.

% =====================================================================
\section{The regime map}
\label{sec:regime}

\begin{table}[t]
\centering
\caption{Per equation: the architecture with the lowest mean relative $L_2$, its
bootstrap interval, the identity floor, whether the best model beats that floor, and the
reference solver's own error over the same horizon. A \textbf{no} in the fourth column is
the important entry: on that equation nothing measured here learned anything useful at
this scale.}
\label{tab:regime}
\small
\begin{tabular}{llccccc}
\toprule
Equation & Best architecture & rel-$L_2$ & 95\% CI & Identity floor & Beats floor?
  & Solver error \\
\midrule
%% Generated by `python -m pinca_jax.paper`. Do not edit.
Advection--diffusion & FNO & 0.00657 & -- & -- & -- & 0.026 \\
Allen--Cahn & FNO & 0.00685 & -- & -- & -- & 0.0008 \\
Cahn--Hilliard & ResNet & 0.0435 & [0.041, 0.046] & 0.412 & yes & 0.023 \\
FitzHugh--Nagumo & NCA & 0.125 & -- & -- & -- & 0.018 \\
Gray--Scott & FNO & 0.674 & -- & -- & -- & 0.011 \\
Heat & Spectral PI-NCA & 0.0417 & [0.039, 0.044] & 0.364 & yes & 0.013 \\
Nagumo & NCA & 0.0733 & -- & -- & -- & 0.001 \\
Navier--Stokes & U-Net & 0.115 & [0.094, 0.14] & 0.156 & yes & 0.01 \\
Shallow water & Multi-channel PI-NCA & 0.0161 & -- & -- & -- & 3.5e-07 \\
Wave & NCA & 0.0521 & -- & -- & -- & 0.0057 \\

\bottomrule
\end{tabular}
\end{table}

Table~\ref{tab:regime} is the headline result, and it should be read in three passes. Left
to right, each row is the decision rule of Section~1 applied to one equation. Down the
``beats floor'' column, any row reading \textbf{no} is an equation on which nothing measured
here learned anything useful, so the winner named in that row is the ordering of noise and
must not be quoted as a result. Down the last column, any row flagged as a non-converging
teacher is an equation whose distillation target is not a converged solution
(Section~\ref{sec:teacher}), which invalidates the row for the same reason by a different
route --- and, as Section~\ref{sec:regime} records, that is not hypothetical: it happened
here, and the affected result was retracted.

\paragraph{Where locality and conservation win.} On smooth local diffusion the best model
is \bestHeat{} at rel-$L_2$ \bestHeatErr{} (CI \bestHeatCI{}), against an identity floor of
\floorHeat{}. Table~\ref{tab:pairedheat} gives the full ranking with paired verdicts, and
the parameter column is where the interesting comparison sits: the conservative local
models reach this accuracy at a small fraction of the FNO's parameter count, and the
physics-free controls at matched parameter count do measurably worse.

\begin{table}[t]
\centering
\caption{Heat: every architecture, ranked, with paired Wilcoxon tests against the
best-performing model on shared initial conditions and Holm--Bonferroni correction.
\emph{tie} means the difference is not resolvable at this sample size.}
\label{tab:pairedheat}
\small
\begin{tabular}{lrcccrl}
\toprule
Architecture & Params & rel-$L_2$ & 95\% CI & Median & $p$ (Holm) & vs.\ best \\
\midrule
%% Generated by `python -m pinca_jax.paper`. Do not edit.
Spectral PI-NCA & 134225 & 0.0417 & [0.039, 0.044] & 0.0418 & -- & \textbf{reference} \\
Multi-channel PI-NCA & 9936 & 0.0426 & [0.038, 0.047] & 0.0396 & 0.63 & tie \\
Bounded PI-NCA & 4576 & 0.0523 & [0.047, 0.058] & 0.0509 & 0.011 & worse \\
PI-NCA (flux) & 4576 & 0.0524 & [0.047, 0.058] & 0.0509 & 0.0092 & worse \\
Bounded multi-scale PI-NCA & 5520 & 0.0614 & [0.057, 0.066] & 0.0583 & 3.1e-05 & worse \\
Multi-scale PI-NCA & 5520 & 0.0614 & [0.057, 0.066] & 0.0583 & 3.1e-05 & worse \\
FNO & 592897 & 0.0634 & [0.06, 0.067] & 0.0618 & 3.1e-05 & worse \\
ResNet & 74336 & 0.0679 & [0.064, 0.072] & 0.0694 & 3.1e-05 & worse \\
U-Net & 265104 & 0.0838 & [0.078, 0.09] & 0.0813 & 3.1e-05 & worse \\
NCA & 6784 & 0.0885 & [0.077, 0.1] & 0.081 & 3.1e-05 & worse \\
ResNet (iso-param) & 5364 & 0.106 & [0.098, 0.11] & 0.109 & 3.1e-05 & worse \\
FNO (iso-param) & 8433 & 0.137 & [0.12, 0.15] & 0.126 & 3.1e-05 & worse \\
U-Net (iso-param) & 7464 & 0.223 & [0.2, 0.24] & 0.22 & 3.1e-05 & worse \\
Identity (floor) & 1 & 0.364 & [0.35, 0.38] & 0.361 & 3.1e-05 & worse \\

\bottomrule
\end{tabular}
\end{table}

\paragraph{Where global mixing wins.} On globally coupled flow the ordering reverses; the
same table for Navier--Stokes is Table~\ref{tab:pairedns}. The mechanism is visible in the
spectral diagnostic (Appendix~\ref{app:spectral}): the local models' residual error
concentrates in the upper half of the spatial spectrum, which is exactly the content that
sharp-interface and turbulent solutions carry.

\begin{table}[t]
\centering
\caption{Navier--Stokes, same protocol as Table~\ref{tab:pairedheat}.}
\label{tab:pairedns}
\small
\begin{tabular}{lrcccrl}
\toprule
Architecture & Params & rel-$L_2$ & 95\% CI & Median & $p$ (Holm) & vs.\ best \\
\midrule
%% Generated by `python -m pinca_jax.paper`. Do not edit.
U-Net & 265104 & 0.115 & [0.094, 0.14] & 0.102 & -- & \textbf{reference} \\
ResNet & 74336 & 0.135 & [0.11, 0.16] & 0.127 & 0.0027 & worse \\
Multi-scale PI-NCA & 5520 & 0.139 & [0.12, 0.16] & 0.136 & 0.0052 & worse \\
Bounded multi-scale PI-NCA & 5520 & 0.14 & [0.12, 0.16] & 0.137 & 0.0052 & worse \\
ResNet (iso-param) & 5364 & 0.145 & [0.12, 0.17] & 0.153 & 0.00076 & worse \\
Multi-channel PI-NCA & 9936 & 0.148 & [0.12, 0.18] & 0.159 & 0.001 & worse \\
Bounded PI-NCA & 4576 & 0.148 & [0.12, 0.18] & 0.153 & 0.00076 & worse \\
PI-NCA (flux) & 4576 & 0.148 & [0.12, 0.18] & 0.152 & 0.00076 & worse \\
Spectral PI-NCA & 134225 & 0.152 & [0.12, 0.18] & 0.155 & 0.00015 & worse \\
FNO (iso-param) & 8433 & 0.152 & [0.12, 0.18] & 0.16 & 0.0013 & worse \\
FNO & 592897 & 0.152 & [0.12, 0.18] & 0.159 & 0.00076 & worse \\
U-Net (iso-param) & 7464 & 0.155 & [0.12, 0.19] & 0.163 & 0.00076 & worse \\
NCA & 6784 & 0.155 & [0.12, 0.19] & 0.164 & 0.00076 & worse \\
Identity (floor) & 1 & 0.156 & [0.12, 0.19] & 0.163 & 0.001 & worse \\

\bottomrule
\end{tabular}
\end{table}

\paragraph{Where the physics-free control wins, and a result we had to retract.} On the
stiff fourth-order equation the best architecture in Table~\ref{tab:regime} is
\bestCahnHilliard{} at \bestCahnHilliardErr{} (CI \bestCahnHilliardCI{}) against a floor of
\floorCahnHilliard{}; the ``beats floor'' entry reads \beatsFloorCahnHilliard{}, and
Table~\ref{tab:pairedch} gives the ranking with paired verdicts. Two things about this row
are worth more than the ranking itself.

First, the winner here is a physics-free convolutional network, with the spectral models
statistically tied to it and the conservative local models measurably worse. That is a
direct counterexample to clause (i) of the hypothesis on this equation, and it is a finding
that exists only because the physics-free controls are in the table; without them the same
data would have shown a conservative PI-NCA ``winning'' a field of PI-NCAs.

Second, an earlier version of this comparison reported that \emph{no} architecture beat the
identity floor here, and wrote that up as a property of the models. It was a property of
the teacher. Cahn--Hilliard is fourth order, so the explicit stability limit is
$\Delta t \le 2/(|\lambda_{\text{lap}}| + \varepsilon^2\lambda_{\text{lap}}^2) = 0.231$,
and the reference implementation shipped $\Delta t = 0.5$. It never visibly exploded only
because the stepper clips the state to $[-1,1]$ every step; remove the clip and the same
configuration reaches NaN. Measured by timestep refinement (Section~\ref{sec:teacher}) it
had an observed order of accuracy of \emph{zero} --- refining $\Delta t$ did not move the
solution at all, because the clip rather than the discretisation was determining the
trajectory. Every architecture ranking distilled from it was ranking models against solver
noise. At a converging timestep the entire picture inverts and every architecture beats the
floor by roughly an order of magnitude. We report this because the failure mode generalises:
the correctness gate had asserted that the JAX solver equalled the PyTorch reference, which
it did, faithfully reproducing the instability. Equality with a reference implementation is
not the same property as being a converged discretisation, and nothing checked the second
one until the teacher-error analysis existed.

\begin{table}[t]
\centering
\caption{Cahn--Hilliard, same protocol as Table~\ref{tab:pairedheat}.}
\label{tab:pairedch}
\small
\begin{tabular}{lrcccrl}
\toprule
Architecture & Params & rel-$L_2$ & 95\% CI & Median & $p$ (Holm) & vs.\ best \\
\midrule
%% Generated by `python -m pinca_jax.paper`. Do not edit.
ResNet & 74336 & 0.0435 & [0.041, 0.046] & 0.0425 & -- & \textbf{reference} \\
Spectral PI-NCA & 134225 & 0.0458 & [0.044, 0.047] & 0.0465 & 0.051 & tie \\
FNO & 592897 & 0.0499 & [0.045, 0.055] & 0.0495 & 0.065 & tie \\
U-Net & 265104 & 0.0501 & [0.046, 0.054] & 0.0491 & 0.00021 & worse \\
Multi-channel PI-NCA & 9936 & 0.0691 & [0.066, 0.072] & 0.0679 & 3.1e-05 & worse \\
ResNet (iso-param) & 5364 & 0.0716 & [0.068, 0.076] & 0.0716 & 3.1e-05 & worse \\
PI-NCA (flux) & 4576 & 0.0831 & [0.08, 0.086] & 0.0825 & 3.1e-05 & worse \\
Bounded PI-NCA & 4576 & 0.0831 & [0.08, 0.086] & 0.0825 & 3.1e-05 & worse \\
NCA & 6784 & 0.0833 & [0.081, 0.086] & 0.0843 & 3.1e-05 & worse \\
Bounded multi-scale PI-NCA & 5520 & 0.0875 & [0.085, 0.09] & 0.0882 & 3.1e-05 & worse \\
Multi-scale PI-NCA & 5520 & 0.0875 & [0.085, 0.09] & 0.0882 & 3.1e-05 & worse \\
FNO (iso-param) & 8433 & 0.152 & [0.14, 0.16] & 0.151 & 3.1e-05 & worse \\
U-Net (iso-param) & 7464 & 0.163 & [0.16, 0.17] & 0.166 & 3.1e-05 & worse \\
Identity (floor) & 1 & 0.412 & [0.4, 0.42] & 0.409 & 3.1e-05 & worse \\

\bottomrule
\end{tabular}
\end{table}

\limitation{The scale reported here is set by the available hardware, and is stated in
Section~\ref{sec:setup} rather than implied. Rankings that are ties at this sample size
may separate at larger grids, longer training and more seeds; rankings that are
significant here may not survive a change of scale. \code{docs/claims\_audit.md} records
which claims the current inventory actually supports, and the two currently unsupported
by the released results are the seed count and the backend.}

% =====================================================================
\section{Do the priors cause the effect? Ablations}
\label{sec:ablate}

A regime map is a correlation until the priors are switched independently at matched
capacity. Each ablation below changes exactly one thing.

\paragraph{A4: conservation, at matched backbone.} Identical perception and processing
width, identical depth; the only difference is whether the output head emits a flux whose
divergence updates the state (mass-conserving by construction on a periodic grid) or a
state increment directly.

\begin{table}[h]
\centering
\caption{A4---conservation on/off at matched backbone width, on heat.}
\label{tab:a4}
\small
\begin{tabular}{lrccc}
\toprule
Head & Params & rel-$L_2$ & 95\% CI & Mass-conservation error \\
\midrule
%% Generated by `python -m pinca_jax.paper`. Do not edit.
flux head & 4576 & 0.0278 & -- & 0.00043 \\
residual head & 4544 & 0.173 & -- & 37 \\

\bottomrule
\end{tabular}
\end{table}

\paragraph{A5: receptive field per step.} Same head, same widths; only the perception
stencil changes, from $3\times3$ to $5\times5$ to dilated multi-scale $(1,2,4)$. This
separates ``locality is the wrong prior'' from ``one cell per step is too slow'', which a
comparison against a global operator alone cannot do.

\begin{table}[h]
\centering
\caption{A5---perception size at matched head and width, on heat.}
\label{tab:a5}
\small
\begin{tabular}{lrccc}
\toprule
Perception & Params & rel-$L_2$ & 95\% CI & Mass-conservation error \\
\midrule
%% Generated by `python -m pinca_jax.paper`. Do not edit.
\multicolumn{5}{l}{\emph{Heat (local diffusion)}} \\
\quad \textbf{$3\times3$} & 4576 & 0.0243 & [0.021, 0.027] & 9.6e-05 \\
\quad $5\times5$ & 5088 & 0.0501 & [0.047, 0.054] & 7.3e-05 \\
\quad dilated $(1,2,4)$ & 9312 & 0.0378 & [0.036, 0.04] & 8.1e-05 \\
\multicolumn{5}{l}{\emph{Navier--Stokes (globally coupled)}} \\
\quad $3\times3$ & 4576 & 0.144 & [0.12, 0.17] & 8.2e-06 \\
\quad $5\times5$ & 5088 & 0.122 & [0.11, 0.14] & 7.4e-06 \\
\quad \textbf{dilated $(1,2,4)$} & 9312 & 0.0961 & [0.084, 0.11] & 6.6e-06 \\

\bottomrule
\end{tabular}
\end{table}

The remaining ablations---iso-parameter spectral mixing (A2), output bounding (A1),
initialisation and pre-seeding protocol (A6)---are in Appendix~\ref{app:ablations}.

% =====================================================================
\section{Bounded and conserving: a defect and a fix}
\label{sec:proj}

Stiff bounded fields need the state clipped to its physical range to remain stable, and
conservative models need total mass restored after that clip. The standard composition is
to clip and then add a uniform offset $\delta = (M_{\text{target}} - \sum u)/N$ to every
cell. That composition is not sound:

\begin{equation}
\text{clip}(u) \in [\ell, h]^N,
\qquad
\text{clip}(u) + \delta \notin [\ell, h]^N \quad\text{whenever } \delta \neq 0,
\end{equation}

because a uniform offset moves every cell, including the ones the clip has just pinned to a
bound. The model is then neither bounded nor, in any useful sense, both.

\emph{How large is this in practice?} We are careful here, because an earlier draft of this
paper quoted a violation of a few percent of cells on Cahn--Hilliard and that figure came
from the non-converging teacher described above. Re-measured against the corrected teacher
(Table~\ref{tab:stability}, out-of-range column) the uniform projection's violation on this
equation is below measurement --- $0\%$ of cells --- because the per-step mass deficit is
small enough that the offset does not push a clipped cell past the bound at this scale. The
defect is therefore a \emph{provable} property of the composition rather than a large
measured effect here: the constructed test in
\code{tests/test\_rigor.py::test\_uniform\_projection\_leaves\_the\_box} exhibits a state on
which it fails, and the violation grows with the size of the deficit the projection has to
absorb, which grows with the timestep, the rollout length, and how hard the clip is binding.
We present the headroom projection as a correctness guarantee that costs nothing, not as a
measured accuracy win, and ablation A7 (Table~\ref{tab:a7}) reports what it actually
changes.

We replace the uniform offset with one proportional to each cell's remaining headroom:

\begin{equation}
u' = u + \operatorname{sign}(d)\,\min\!\left(\frac{|d|}{C},\,1\right) \cdot C \cdot
     \frac{r}{C},
\qquad
r = \begin{cases} h - u & d \geq 0\\ u - \ell & d < 0\end{cases},
\qquad
C = \textstyle\sum r,
\end{equation}

where $d = M_{\text{target}} - \sum u$. Mass is restored exactly whenever the target is
feasible for the box, no cell can cross the bound because the total increment is capped by
the total headroom by construction, the map is differentiable, and it is a single
vectorised pass. When the target is genuinely infeasible---$M_{\text{target}} \notin
[N\ell, Nh]$---the field saturates at the bound and the residual mass error is
\emph{reported} rather than paid for by violating the bound. Ablation A7 measures the cost
of the naive choice rather than asserting it.

\begin{table}[h]
\centering
\caption{A7---how mass is restored after the bound clip, on Cahn--Hilliard. Same backbone,
same bounds, one line different.}
\label{tab:a7}
\small
\begin{tabular}{lrccc}
\toprule
Projection & Params & rel-$L_2$ & 95\% CI & Mass-conservation error \\
\midrule
%% Generated by `python -m pinca_jax.paper`. Do not edit.
\multicolumn{1}{l}{\emph{ablation A7 on cahn\_hilliard: not yet measured. Run \texttt{python -m pinca\_jax.bench\_all --group ablation --seeds 5}.}}

\bottomrule
\end{tabular}
\end{table}

% =====================================================================
\section{Stability: failures counted, not clamped}
\label{sec:stability}

Accuracy tables in this paper clamp rollouts to the reference solver's measured physical
range, so a diverging model cannot emit a meaningless metric---a negative PSNR is only
ever the statement that mean squared error exceeded the signal range. That guard is
correct for accuracy and wrong for stability, because it converts a blow-up into a
saturated field and a merely-poor score. Every measurement in this section therefore runs
with the guard \emph{disabled}.

\begin{table}[h]
\centering
\caption{Stability under stress on Cahn--Hilliard, divergence guard disabled. A rollout is
failed once it is non-finite or exceeds ten times the reference solver's physical range.
Perturbation amplification is the growth of a $10^{-3}$ relative perturbation over the
rollout.}
\label{tab:stability}
\small
\begin{tabular}{lrcccc}
\toprule
Architecture & Params & Failure rate & Median survival & Out-of-range cells
  & Perturbation amp. \\
\midrule
%% Generated by `python -m pinca_jax.paper`. Do not edit.
NCA & 6784 & 100\% & 58/96 & 0.00\% & 1.82$\times$ \\
PI-NCA (flux) & 4576 & 100\% & 62/96 & 0.00\% & 1.8$\times$ \\
Multi-scale PI-NCA & 5520 & 100\% & 60/96 & 0.00\% & 1.8$\times$ \\
Bounded multi-scale PI-NCA & 5520 & 0\% & 96/96 & 0.00\% & 1.8$\times$ \\
uniform projection & 5520 & 0\% & 96/96 & 0.00\% & 1.8$\times$ \\
headroom projection & 5520 & 0\% & 96/96 & 0.00\% & 1.8$\times$ \\
FNO & 592897 & 100\% & 73/96 & 0.00\% & 1.79$\times$ \\
ResNet & 74336 & 100\% & 66/96 & 0.00\% & 1.79$\times$ \\
U-Net & 265104 & 100\% & 69/96 & 0.00\% & 1.78$\times$ \\
Identity (floor) & 1 & 0\% & 96/96 & 0.00\% & 1$\times$ \\

\bottomrule
\end{tabular}
\end{table}

The result is categorical rather than marginal: at \stabHorizon{} steps the unconstrained
local model fails on \stabNCAFail\% of initial conditions while the bounded conservative
variant fails on \stabBoundedFail\%. Note that this ordering is the opposite of what
Table~\ref{tab:pairedch} shows on accuracy at the shorter horizon, which is the point.
Stability and accuracy are different axes and a paper that reports only the second will
recommend the wrong model for any application that runs long rollouts.

\limitation{Failure is defined by a finite threshold (ten times the reference range) and a
finite horizon. A model that degrades gracefully without crossing that threshold is
recorded as surviving; the accompanying accuracy and out-of-range columns are what
distinguish surviving from being useful.}

% =====================================================================
\section{Is the distillation target meaningful?}
\label{sec:teacher}

The emulators are trained to imitate a numerical solver and then scored against that same
solver, so a low relative $L_2$ could mean accurate physics or a faithful copy of a poor
approximation. We measure the difference.

For the linear constant-coefficient members of the suite the periodic problem is diagonal
in Fourier space, so the continuum solution is available in closed form and the solver's
absolute error is directly computable. We report it split. Against the exactly-integrated
\emph{semi-discrete} system (the 5-point Laplacian's own eigenvalues) the gap is the
forward-Euler temporal error; against the continuum solution it is the total, and the
difference between them is the stencil's spatial truncation error. The split matters:
refining $\Delta t$ converges the solver to the semi-discrete system, not to the
continuum, so the spatial error is a floor and reporting only the continuum gap makes that
floor look like a failure to converge. For the nonlinear members no closed form exists, so
the teacher is refined in time and the Cauchy differences between refinements bound its
own error, with the observed order of accuracy reported alongside.

\begin{table}[h]
\centering
\caption{The reference solver's own error over the evaluation horizon, and its cost.
\emph{vs analytic}: closed-form Fourier solution. \emph{self convergence}: relative $L_2$
between the production $\Delta t$ and a $\Delta t/8$ reference over the same physical
time.}
\label{tab:teacher}
\small
\begin{tabular}{llccccrc}
\toprule
Equation & Estimate & Total & Spatial & Temporal & Order & Solver ms/step & Speedup \\
\midrule
%% Generated by `python -m pinca_jax.paper`. Do not edit.
Heat & vs analytic & 0.013 & 0.02 & 0.0072 & -- & 0.040 & 0.05$\times$ \\
Wave & vs analytic & 0.0057 & 0.0059 & 0.0016 & -- & 0.025 & 0.04$\times$ \\
Advection--diffusion & vs analytic & 0.026 & 0.027 & 0.0018 & -- & 0.078 & 0.12$\times$ \\
Allen--Cahn & self convergence & 0.0008 & -- & -- & 1.04 & 0.035 & 0.04$\times$ \\
Gray--Scott & self convergence & 0.011 & -- & -- & 1.02 & 0.108 & 0.12$\times$ \\
Shallow water & self convergence & 3.5e-07 & -- & -- & -0.56 & 0.165 & 0.31$\times$ \\
Cahn--Hilliard & self convergence & 0.023 & -- & -- & 0.90 & 0.090 & 0.12$\times$ \\
FitzHugh--Nagumo & self convergence & 0.018 & -- & -- & 1.00 & 0.053 & 0.07$\times$ \\
Nagumo & self convergence & 0.001 & -- & -- & 0.99 & 0.015 & 0.02$\times$ \\
Navier--Stokes & self convergence & 0.01 & -- & -- & 1.15 & 0.249 & 0.31$\times$ \\

\bottomrule
\end{tabular}
\end{table}

On heat the total solver error is \teacherHeatTotal{}, of which \teacherHeatSpatial{} is
spatial and \teacherHeatTemporal{} temporal. Emulator errors on the same equation
(Table~\ref{tab:pairedheat}) sit above that, which is the case we want: the models are
limited by learning, not by the target, so comparisons between them are measuring
architecture rather than teacher noise. Where an emulator's error approaches the teacher's,
we say so and stop ranking.

% =====================================================================
\section{Is any of this worth it?}
\label{sec:worth}

Two comparisons that are usually omitted, both of which cut against the surrogates.

\paragraph{Against PINNs, on a matched task.} A PINN solves one initial-value problem per
training run; an emulator trains once and applies to any initial condition. Comparing them
on a single accuracy column compares different quantities. We instead fix the task---given
$K$ initial conditions from the same distribution, produce the solution at time $T$ for all
of them---and report both cost structures.

\begin{table}[h]
\centering
\caption{Matched comparison on heat over \matchedK{} held-out initial conditions: same
equation, same initial conditions, same horizon, same reference, same metric.}
\label{tab:matched}
\small
\begin{tabular}{lccrrrr}
\toprule
Method & rel-$L_2$ & 95\% CI & Params & Fixed cost (s) & Per-IC (s) & Total (s) \\
\midrule
%% Generated by `python -m pinca_jax.paper`. Do not edit.
Emulator (Multi-scale PI-NCA) & 0.161 & [0.15, 0.18] & 5520 & 2.8 & 0.0584 & 3.0 \\
PINN (one run per IC) & 0.318 & [0.29, 0.35] & 14209 & 0.0 & 5.9541 & 23.8 \\
Numerical solver (reference) & 0 & -- & -- & 0 & 0.0000 & 0.000 \\

\bottomrule
\end{tabular}
\end{table}

The number a practitioner needs from this is the crossover: at $K = \matchedCrossover{}$
initial conditions the emulator's one-off training is amortised against per-instance PINN
solving. Below it the PINN is cheaper in total wall-clock; above it the emulator is. A
single-instance accuracy comparison between the two paradigms is therefore not
informative on its own, in either direction.

\paragraph{Against the solver.} The solver row of Table~\ref{tab:matched} is the one that
matters. At the grid sizes used here the reference solver is roughly
\matchedSolverSpeedup$\times$ faster per initial condition than the trained emulator, and
exact by the definition of the comparison. \textbf{At this scale no surrogate in this paper
is a deployment case.} Learned surrogates pay off where the solver does not fit---much
larger grids, much longer horizons, stiff timestep restrictions, or the need to
differentiate through an entire rollout---and none of those conditions holds on a small
periodic grid with an explicit stepper. We state this because the alternative is to let a
favourable architecture ranking imply a practical claim the measurements do not support.

% =====================================================================
\section{Generalisation beyond the training distribution}
\label{sec:ood}

Random draws from the training initial-condition generator are not a test set; they are
more training data with a different seed. We define held-out axes before training and
evaluate on each: amplitude, structure count, spatial length scale, added broadband
spectral content, PDE coefficient, and horizon. A random torus translation is included as
a \emph{negative control}: every architecture here is translation-equivariant by
construction, so a gap on that axis would indicate a bug in the evaluation rather than a
property of a model.

\begin{table}[h]
\centering
\caption{Out-of-distribution degradation on heat, as the ratio of the shifted-axis error
to the in-distribution error. $1.00\times$ is perfect transfer. Bold entries have at least
one initial condition whose relative $L_2$ exceeds 1, i.e.\ worse than predicting no
change.}
\label{tab:ood}
\small
%% Generated by `python -m pinca_jax.paper`. Do not edit.
\begin{tabular}{lcccccc}
\toprule
Held-out axis & NCA & PI-NCA (flux) & Multi-scale PI-NCA & FNO & ResNet & U-Net \\
\midrule
in\_dist & 1.00$\times$ & 1.00$\times$ & 1.00$\times$ & 1.00$\times$ & 1.00$\times$ & 1.00$\times$ \\
phase & 1.00$\times$ & 1.00$\times$ & 1.00$\times$ & 1.00$\times$ & 1.00$\times$ & 1.00$\times$ \\
amp\_half & 1.06$\times$ & 1.00$\times$ & 0.99$\times$ & 1.24$\times$ & 0.93$\times$ & 0.95$\times$ \\
amp\_2x & 0.95$\times$ & 1.32$\times$ & 1.38$\times$ & 1.58$\times$ & 1.47$\times$ & 1.43$\times$ \\
amp\_4x & 1.60$\times$ & 3.35$\times$ & 3.46$\times$ & 3.31$\times$ & 3.61$\times$ & 3.06$\times$ \\
blobs\_few & 1.14$\times$ & 0.98$\times$ & 0.99$\times$ & 1.33$\times$ & 1.01$\times$ & 0.98$\times$ \\
blobs\_many & 0.75$\times$ & 0.90$\times$ & 0.92$\times$ & 1.33$\times$ & 0.97$\times$ & 1.02$\times$ \\
scale\_narrow & \textbf{7.45$\times$} & 6.90$\times$ & 9.06$\times$ & 6.57$\times$ & \textbf{12.00$\times$} & \textbf{10.10$\times$} \\
scale\_wide & 0.62$\times$ & 0.24$\times$ & 0.30$\times$ & 0.93$\times$ & 0.33$\times$ & 0.31$\times$ \\
spectrum\_rough & 1.21$\times$ & 1.60$\times$ & 1.52$\times$ & 1.41$\times$ & 1.63$\times$ & 1.42$\times$ \\
coeff\_half & 0.59$\times$ & 1.43$\times$ & 1.58$\times$ & 1.18$\times$ & 1.41$\times$ & 0.97$\times$ \\
coeff\_2x & 2.44$\times$ & 3.76$\times$ & 3.48$\times$ & 3.24$\times$ & 3.84$\times$ & 3.29$\times$ \\
horizon\_2x & 1.26$\times$ & 1.24$\times$ & 1.19$\times$ & 1.09$\times$ & 1.04$\times$ & 0.96$\times$ \\
horizon\_4x & 1.25$\times$ & 1.54$\times$ & 1.93$\times$ & 1.40$\times$ & 1.17$\times$ & 0.92$\times$ \\
\bottomrule
\end{tabular}

\end{table}

Resolution transfer is the seventh axis and is reported separately in
Appendix~\ref{app:resolution}, because the convolutional and spectral models are
size-agnostic and can be evaluated at grids they were never trained on.

% =====================================================================
\section{Efficiency, within a budget rather than across one}
\label{sec:efficiency}

Comparing a spectral operator with $5.9\times10^5$ parameters against a cellular automaton
with $5\times10^3$ and calling the difference an architecture result smuggles a capacity
result into an inductive-bias claim. We therefore group by parameter budget and compare
within a class first.

\begin{table}[h]
\centering
\caption{Heat, grouped by parameter budget. Within a class the comparison isolates the
prior; across classes it reports what the extra parameters bought.}
\label{tab:efficiency}
\small
\begin{tabular}{lrccrr}
\toprule
Architecture & Params & rel-$L_2$ & 95\% CI & Train (s) & Infer (ms/step) \\
\midrule
%% Generated by `python -m pinca_jax.paper`. Do not edit.
\multicolumn{6}{l}{\emph{floor budget}} \\
\quad Identity (floor) & 1 & 0.364 & [0.35, 0.38] & 1 & 0.011 \\
\multicolumn{6}{l}{\emph{small budget}} \\
\quad Multi-channel PI-NCA & 9936 & 0.0426 & [0.038, 0.047] & 3 & 0.938 \\
\quad Bounded PI-NCA & 4576 & 0.0523 & [0.047, 0.058] & 2 & 0.315 \\
\quad PI-NCA (flux) & 4576 & 0.0524 & [0.047, 0.058] & 2 & 0.304 \\
\quad Bounded multi-scale PI-NCA & 5520 & 0.0614 & [0.057, 0.066] & 2 & 0.524 \\
\quad Multi-scale PI-NCA & 5520 & 0.0614 & [0.057, 0.066] & 2 & 0.501 \\
\quad NCA & 6784 & 0.0885 & [0.077, 0.1] & 2 & 0.283 \\
\quad ResNet (iso-param) & 5364 & 0.106 & [0.098, 0.11] & 2 & 1.406 \\
\quad FNO (iso-param) & 8433 & 0.137 & [0.12, 0.15] & 2 & 0.387 \\
\quad U-Net (iso-param) & 7464 & 0.223 & [0.2, 0.24] & 3 & 1.159 \\
\multicolumn{6}{l}{\emph{medium budget}} \\
\quad ResNet & 74336 & 0.0679 & [0.064, 0.072] & 17 & 2.788 \\
\multicolumn{6}{l}{\emph{large budget}} \\
\quad Spectral PI-NCA & 134225 & 0.0417 & [0.039, 0.044] & 3 & 1.259 \\
\quad FNO & 592897 & 0.0634 & [0.06, 0.067] & 4 & 2.672 \\
\quad U-Net & 265104 & 0.0838 & [0.078, 0.09] & 13 & 4.426 \\

\bottomrule
\end{tabular}
\end{table}

The within-class comparison is the one that speaks to the hypothesis: at a matched
parameter budget the conservative local models are compared against a plain CNN, a U-Net
and a spectral operator reduced to the same budget, so any gap there is attributable to the
prior rather than to capacity. The across-class rows then answer the separate question of
what a two-order-of-magnitude increase in parameters actually buys on this equation, which
is not the same question and should not be reported as though it were.

% =====================================================================
\section{Does the ranking survive a change of scale?}
\label{sec:scaling}

Every table above is measured at one grid, one training horizon and one epoch budget. A
ranking read off a single operating point is only quotable if it is stable under scale,
and there is no reason to assume that: a spectral operator retains a fixed number of
Fourier modes and should improve with resolution differently from a local rule whose
receptive field is fixed in cells, and a model with two orders of magnitude more
parameters should respond differently to a larger training budget. We therefore sweep each
axis with everything else held fixed and report Kendall's $\tau$ between the ordering at
the smallest and the largest setting.

\begin{table}[h]
\centering
\caption{Rank stability along each scale axis on heat. $\tau = 1$ means the ordering is
unchanged; anything less means the paper's ranking is conditional on its operating point.}
\label{tab:scaling}
\small
%% Generated by `python -m pinca_jax.paper`. Do not edit.
\begin{tabular}{lllcl}
\toprule
Axis & Swept & Best at each end & Kendall $\tau$ & Verdict \\
\midrule
grid & 16--24 & PI-NCA (flux) $\to$ FNO & 0.67 & \textbf{ordering changes} \\
rollout & 4--8 & PI-NCA (flux) $\to$ FNO & 0.67 & \textbf{ordering changes} \\
epochs & 60--150 & PI-NCA (flux) $\to$ PI-NCA (flux) & 1.00 & stable \\
\bottomrule
\end{tabular}

\end{table}

This is the section that most constrains how the rest of the paper may be read. Where
$\tau < 1$ the ordering between architectures changes with the axis, and the headline
ranking must be quoted together with the grid and horizon it was measured at rather than
as a property of the architectures. Currently the ordering is not stable along
\scalingUnstableAxes{}, which is exactly what one would predict from the mechanism in
Appendix~\ref{app:spectral}: the local models' fixed per-step receptive field means the
number of steps required for information to cross the domain grows with the grid, so their
advantage on a small domain is partly an artefact of the domain being small.

\limitation{Kendall's $\tau$ compares only the endpoints of each sweep and treats all rank
inversions as equal, including inversions between architectures whose intervals overlap.
An ordering flagged unstable here may be a flip between two models the paired tests
already call a tie; the per-setting errors in \code{results/scaling\_*.md} are what
distinguish the two cases.}

% =====================================================================
\section{Three dimensions}
\label{sec:threed}

The suite is repeated in 3-D on \NpdesThreeD{} families with the same harness, metrics and
protocol (NDHWC arrays, 7-point Laplacian, dimension-appropriate stability limits).

\begin{table}[h]
\centering
\caption{3-D results: best architecture per equation, with the identity floor.}
\label{tab:threed}
\small
\begin{tabular}{llcc}
\toprule
Equation & Best architecture & rel-$L_2$ & Identity floor \\
\midrule
%% Generated by `python -m pinca_jax.paper`. Do not edit.
Advection--diffusion & FNO & 0.00844 & -- \\
Allen--Cahn & FNO & 0.0119 & -- \\
FitzHugh--Nagumo & FNO & 0.163 & -- \\
Gray--Scott & NCA & 0.418 & -- \\
Heat & PI-NCA (flux) & 0.0467 & -- \\
Nagumo & NCA & 0.0407 & -- \\

\bottomrule
\end{tabular}
\end{table}

\limitation{We deliberately do \emph{not} describe the agreement between
Tables~\ref{tab:regime} and \ref{tab:threed} as evidence that the regime dependence is
``dimension-independent''. The 3-D suite is smaller, at a single reduced resolution, and
the two are not independent samples of the same population. The correct statement is that
the 3-D orderings are consistent with the 2-D ones on the equations measured in both.}

% =====================================================================
\section{Limitations}
\label{sec:limits}

Stated as a list, because each of these bounds a claim made above.

\begin{enumerate}
\item \textbf{Scale.} Grids \Grids{} and \MinSeeds--\MaxSeeds{} seeds per table on backend
  \Backends{}. Ties may separate and significant orderings may not survive at larger
  scale. \code{docs/claims\_audit.md} records exactly which pre-registered claims the
  released inventory supports; the seed count and the backend are currently among those it
  does not.
\item \textbf{Distillation, not discovery.} Every emulator is trained against a solver, so
  none of them can be more accurate than that solver plus its own error. The teacher's
  error is measured (Section~\ref{sec:teacher}) so the ceiling is known, not assumed.
\item \textbf{Periodic regular grids only.} No irregular geometry, no non-periodic
  boundaries. Graph- and mesh-based surrogates are therefore not compared, and no claim
  here extends to them.
\item \textbf{PINO is not measured.} It is the most obviously missing baseline on the
  physics-informed axis. The PINN and DeepONet baselines here are heat-only and reported
  separately.
\item \textbf{One initial-condition generator per equation.} The out-of-distribution axes
  (Section~\ref{sec:ood}) perturb it along named directions but do not replace it with an
  independently designed family.
\item \textbf{No shared public benchmark.} The solvers here are our own, so the numbers are
  not comparable to published PDEBench or APEBench results.
\item \textbf{The surrogates lose to the solver at this scale} (Section~\ref{sec:worth}).
  The regimes where they would not are precisely the ones this hardware cannot reach.
\end{enumerate}

% =====================================================================
\section{Conclusion}

We set out to test whether the best structural prior for a neural PDE surrogate is
determined by the equation rather than by capacity. The evidence supports a qualified
version: the ordering of architectures does change with the equation's structure in the
directions the hypothesis predicts, and it does so at parameter counts that differ by two
orders of magnitude, so capacity is not the explanation. But the controls change the
picture in ways a comparison without them would have missed---on at least one equation
nothing measured beats a do-nothing baseline, physics-free CNNs are competitive in
regimes where a physics prior was expected to dominate, accuracy and long-horizon
stability rank models differently, and at the scales measured here the numerical solver
beats every surrogate on both accuracy and cost. The useful output of this work is
therefore not an architecture but a protocol: the floor, the physics-free controls, the
paired statistics, the guard-off stability measurement, and the teacher-error accounting
are what turn a table of numbers into a claim that can be checked, and they are what we
release.

\subsubsection*{Reproducibility statement}
All code, raw results and generated tables are in the repository. \code{python -m pytest
tests/} runs \Ntests{} correctness tests asserting the ported solvers and models equal
verbatim references to tolerance. \code{python -m pinca\_jax.runner} runs the entire
pipeline end to end and is resumable per matrix cell. \code{python -m pinca\_jax.claims}
regenerates the audit of every claim in this paper against the raw results; \code{python
-m pinca\_jax.paper} regenerates every table and every quoted number in this document;
\code{python -m pinca\_jax.bib} re-verifies every citation against arXiv. Each results
file records the JAX version, backend, devices, peak memory and exact configuration.

\bibliographystyle{plainnat}
\bibliography{refs}

\appendix
% Appendices. Engineering detail lives here; the main paper carries the evidence.
% Generated tables are pulled in by \input, never retyped.

% =====================================================================
\section{Related work: full positioning table}\label{app:related}

Locality: does the update read a bounded neighbourhood per step? Conservation: enforced
structurally (\textbf{S}), by projection (\textbf{P}), by penalty (\textbf{L}), or not at
all. The full discussion, including the scope limits this work does not attempt, is in
\code{docs/related\_work.md}.

\begin{table}[h]
\centering
\small
\begin{tabular}{llllll}
\toprule
Work & Locality & Conservation & Supervision & Suite & Dims \\
\midrule
Growing NCA (Mordvintsev et al., 2020) & local CA rule & -- & target image & n/a & 2-D \\
NCA time-stepping \citep{saha2026nca} & local CA rule & -- & solver traj. & 5 PDEs & 1/2-D \\
Spatio-temporal NCA \citep{richardson2023nca} & local CA rule & symmetry & traj., images & Turing & 2-D \\
NoiseNCA \citep{pajouheshgar2024noisenca} & local CA rule & -- & texture & n/a & 2-D \\
FINN \citep{praditia2021finn,karlbauer2022composing} & FV stencil & \textbf{S} & data + physics & transport & 1/2-D \\
Neural Conservation Laws \citep{richterpowell2022ncl} & global forms & \textbf{S} & data & fluids & 2/3-D \\
PDE-Net \citep{long2018pdenet,long2019pdenet2} & local filters & -- & trajectories & canonical & 2-D \\
FNO \citep{li2021fno} & \textbf{global} & -- & input--output & canonical & 2/3-D \\
PINO \citep{li2021pino} & global & \textbf{L} & data + residual & canonical & 2-D \\
DeepONet \citep{lu2021deeponet} & global & -- & operator samples & canonical & 1/2-D \\
MP-PDE \citep{brandstetter2022mppde} & message passing & -- & trajectories & canonical & 1/2-D \\
MeshGraphNets \citep{pfaff2021meshgraphnets} & mesh & -- & trajectories & meshes & 2/3-D \\
GNS \citep{sanchez2020gns} & graph & -- & trajectories & particles & 2/3-D \\
Clifford layers \citep{brandstetter2023clifford} & local & geometric & trajectories & fluid/EM & 2/3-D \\
PDE-Refiner \citep{lippe2023refiner} & global & -- & trajectories & canonical & 1/2-D \\
PINN \citep{raissi2019pinn} & coordinate MLP & \textbf{L} & residual only & per-IVP & any \\
PDEBench \citep{takamoto2022pdebench} & n/a & n/a & benchmark & broad & 1--3D \\
APEBench \citep{koehler2024apebench} & n/a & n/a & benchmark & broad & 1--3D \\
\midrule
\textbf{This work} & local CA rule & \textbf{S} flux, \textbf{P} headroom & distillation
  & \NpdesTwoD{}\,+\,\NpdesThreeD{} & 2/3-D \\
\bottomrule
\end{tabular}
\end{table}

% =====================================================================
\section{PDE suite: equations, parameters, stability limits}\label{app:suite}

All equations are posed on a periodic grid with unit spacing, so the domain is $[0,N)^d$
and the physical wavenumber of index $k$ is $2\pi k/N$. Spatial derivatives use the
5-point Laplacian and central differences in 2-D and the 7-point Laplacian in 3-D; time
integration is forward Euler except for shallow water (RK4) and Navier--Stokes
(pseudo-spectral vorticity form, integrating factor for viscosity, FFT Poisson solve for
the stream function).

Two parameter choices deserve stating rather than burying. Gray--Scott ships in the source
notebook with $\Delta t = 2.0$, which exceeds the explicit-diffusion stability bound
$\Delta t \le \Delta x^2 / (4 D_u) = 1.25$ and diverges on sharp initial conditions by
step 12. The verbatim parameters are retained in the registry so the migration tests can
assert equality against the original, and a separate stable configuration is used for
every experiment ($\Delta t = 1.0$; $0.5$ in 3-D, where the bound is $\Delta x^2/(6D_u)$).
That divergence is itself asserted by a test, so the finding cannot silently regress.

Initial conditions are per-equation: periodic sums of Gaussian blobs (heat,
advection--diffusion, wave, shallow water, Navier--Stokes), smoothed noise (Cahn--Hilliard,
Nagumo), sign fields with noise (Allen--Cahn), random square seeds (Gray--Scott), and small
perturbations of the rest state (FitzHugh--Nagumo). Full generators are in
\code{src/pinca\_jax/ic.py}. The held-out variants used in Section~\ref{sec:ood} live in a
separate module, \code{src/pinca\_jax/ood.py}, so the correctness-gated generators stay
untouched.

% =====================================================================
\section{Architectures: update equations}\label{app:arch}

Let $u \in \mathbb{R}^{H\times W\times C}$ be the state, $P_\theta$ a perception stage
(circularly padded convolution) and $F_\theta$ a pointwise processing stage.

\paragraph{Unconstrained NCA.} $u' = u + W_{\text{out}} F_\theta(P_\theta(u))$ with
$W_{\text{out}}$ zero-initialised, so the model begins as the identity. No conservation.

\paragraph{Conservative flux PI-NCA.} The head emits a flux field
$f \in \mathbb{R}^{H\times W\times 2C}$ and the update is its discrete divergence,
\begin{equation}
u'_c = u_c + \big(\mathrm{roll}(f^x_c, 1, W) - f^x_c\big)
           + \big(\mathrm{roll}(f^y_c, 1, H) - f^y_c\big),
\end{equation}
which telescopes to exactly zero when summed over a periodic grid, so $\sum u'_c = \sum
u_c$ per channel up to floating-point error. \emph{Per channel} is the operative phrase: a
single lumped total would let one field's deficit be paid out of another's, which is wrong
for any multi-field system.

\paragraph{Multi-scale variant.} The perception stage becomes concatenated dilated
convolutions at dilations $(1,2,4)$, widening the per-step receptive field from $\pm1$ to
$\pm4$ cells without an FFT and without a large parameter increase.

\paragraph{Spectral variant.} A truncated-Fourier global stream is added to the local
conservative stream and the sum is mass-projected. This is the ``global operator plus
local conservation'' hypothesis, tested directly rather than argued.

\paragraph{Bounded variants.} As above, followed by a clip to the equation's measured
physical range and then the mass projection of Section~\ref{sec:proj}.

\paragraph{Physics-free controls.} A residual CNN (circular padding, no resolution change,
receptive field $1 + 4\,\text{depth}$ cells per step) and a U-Net (mean-pool down,
nearest-neighbour up, circular padding, skip connections). Both carry zero-initialised
heads, so they also begin as the identity, and both are additionally instantiated at
parameter counts matched to the NCAs.

\paragraph{Bounds are measured, not assumed.} The bounded variants take their range from a
short reference-solver rollout rather than a hardcoded $[-1,1]$. Hardcoding is correct for
Cahn--Hilliard and Allen--Cahn and destroys a field whose amplitudes run 5--10, which is
why bounding could previously be evaluated only on the two equations it was tuned for.

% =====================================================================
\section{Remaining ablations}\label{app:ablations}

\textbf{A1 (output bounding)} switches the clip on and off on the stiff bounded equations.
\textbf{A2 (iso-parameter spectral mixing)} compares the full FNO against one reduced to
the NCA parameter budget, separating ``spectral mixing helps'' from ``more parameters
help''; both appear in the main tables, as \code{fno} and \code{fno\_small}. \textbf{A6
(initialisation and pre-seeding)} varies the training protocol: He initialisation,
zero-initialised heads, learning-rate warm-up, and pre-evolving the training initial
conditions so the model sees developed states rather than only early ones. Pre-seeding
helps on most equations and measurably hurts on Cahn--Hilliard, whose solutions coarsen
rapidly from fresh initial conditions so that developed states mismatch the evaluation
distribution. It is therefore disabled for that equation, and the exception is recorded in
the configuration rather than applied silently.

% =====================================================================
\section{Spectral error diagnostic}\label{app:spectral}

For each model we compute the fraction of squared error energy in the upper half of the
2-D spatial spectrum. The local models concentrate their residual error at fine scales
while the spectral operator's error is closer to broadband-flat. This is the mechanism
behind the regime split in Table~\ref{tab:regime}: it predicts an advantage for the global
operator on exactly the equations whose solutions carry sharp interfaces, and it does so
from a quantity independent of the ranking it explains.

% =====================================================================
\section{Resolution transfer}\label{app:resolution}

Convolutional NCAs and spectral operators are both size-agnostic, so a model trained at one
grid can be evaluated at another with no retraining. We train at each grid in the study and
evaluate at every grid, giving a transfer matrix whose diagonal is same-grid accuracy and
whose off-diagonal entries are zero-shot coarse-to-fine (super-resolution) and
fine-to-coarse transfer. Raw tables are in \code{results/bench\_resolution\_*.md}.

% =====================================================================
\section{Migration and the correctness gate}\label{app:migration}

Every solver and model in this study was migrated from a PyTorch original and asserted
equal to a verbatim reference to tolerance \emph{before} any experiment was run. The gate
is \Ntests{} tests, and it is the reason a surprising result in this paper can be
attributed to the science rather than to a bug.

It earned its keep. The periodic Laplacian originally differenced the last two axes, which
is correct for the NCHW layout the unit test used and wrong for the NHWC arrays the
training path actually passes, so the teacher was a degenerate one-dimensional Laplacian
while the test still passed. The fix was to difference explicit spatial axes and to add an
isotropy regression test exercising the real data path. We record it because the lesson
generalises: a correctness gate that does not test the layout the pipeline uses can pass
while the pipeline is wrong.

The original PyTorch script is retained at the repository root for provenance only, with
its defects documented at line level in \code{docs/legacy\_pytorch.md} --- most
consequentially a plotted training curve that can be identically zero, because the
accumulated loss is read after the truncation branch has reset it, and a validation loop
that omits the conservation projection which training and the final test both apply, so it
scores a different model. No number in this paper comes from running it.


\end{document}
